% `tex.web` §1073, §1084, §1078: how much of a leader is one command.
%
%     any_mode(leader_ship): scan_box(leader_flag-a_leaders+cur_chr);
%
% §1084's `scan_box` reads the box — or the rule, which is the one place a
% rule is legal in restricted horizontal mode — without going back through
% `main_control`. §1078's "Append a new leader node" then reads the glue and
% calls `append_glue` directly. So `\cleaders\hbox{A}\hskip 20pt` is *one*
% traced command, not three; we traced `{\leaders}` `{\hbox}` `{\hskip}` where
% Knuth's transcript has `{\cleaders}` alone.
%
% And it is `{\cleaders}`: §1078's three differ only in how the box is spaced
% along the glue, but they are three `chr` values in TeX, so `\meaning` and the
% trace name each. We had collapsed them onto one primitive, which made
% `\meaning\xleaders` say `\leaders` — including through a `\let` alias, where
% TeX still prints the primitive it came from.
\catcode`\{=1 \catcode`\}=2
\tracingonline=1
\tracingcommands=1
\let\zz=\cleaders
\message{[\meaning\cleaders][\meaning\xleaders][\meaning\zz][\meaning\leaders]}

\setbox0\hbox{\leaders\hbox{A}\hskip 20pt
  \cleaders\vrule width 2pt\hskip 20pt
  \xleaders\copy0\hskip 20pt
  \zz\hbox{B}\hskip 20pt}

% In vertical mode the glue is a `\vskip`, and the box is still scanned by the
% leader rather than by the mode.
\leaders\hbox{C}\vskip 20pt
\end
